\section{Resultados}
\label{sec:resultados}

\subsection{El pasillo donde todos los mecanismos son idénticos}
\label{sec:pasillo}

En el banco de réplica, recortar el catálogo de nueve clases a seis sube el F1 macro alineado de
\num{0.4895} a \num{0.7342}, una mejora aparente de \num{0.2447}. En todo ese tramo las cuatro
series —retirada por F1, retirada por soporte, rechazo por confianza y el control que no aplica
mecanismo alguno— devuelven el \emph{mismo} valor, y la cobertura se queda en \num{1.000}.

La razón es aritmética y es lo que hace concluyente la observación. Las tres clases que se retiran
primero tienen dos, siete y \num{289} parcelas de \num{60000}; el argmax sin restringir
prácticamente nunca las predice, así que quitarlas de la leyenda no quita ninguna parcela del mapa
entregado. Ningún mecanismo hace nada. La métrica sube casi un cuarto de punto por su cuenta.

\begin{table}[tbp]
\caption{Descomposición de la mejora aparente en el banco de réplica, universo de nueve clases. La
columna «denominador» es el control de la sección~\ref{sec:control}: el predictor intacto puntuado
sobre la misma leyenda, entregándolo todo. $K = 5$ es el criterio principal preregistrado.}
\label{tab:descomposicion}
\centering
\begin{tabular}{rrrrrr}
\toprule
$K$ & F1 alineado & mejora aparente & denominador & mecanismo & cobertura \\
\midrule
9 & 0,4895 & ---     & ---              & ---            & 1,000 \\
8 & 0,5507 & +0,0612 & +0,0612 (100,0\%) & +0,0000 (0,0\%) & 1,000 \\
7 & 0,6293 & +0,1398 & +0,1398 (100,0\%) & +0,0000 (0,0\%) & 1,000 \\
6 & 0,7342 & +0,2447 & +0,2447 (100,0\%) & +0,0000 (0,0\%) & 1,000 \\
\textbf{5} & \textbf{0,8055} & \textbf{+0,3160} & \textbf{+0,3004 (95,1\%)} & \textbf{+0,0156 (4,9\%)} & \textbf{0,859} \\
4 & 0,8579 & +0,3684 & +0,3352 (91,0\%) & +0,0333 (9,0\%) & 0,806 \\
3 & 0,9126 & +0,4232 & +0,3826 (90,4\%) & +0,0406 (9,6\%) & 0,444 \\
\bottomrule
\end{tabular}
\end{table}

\subsection{La descomposición se transporta}
\label{sec:transporte}

El mismo protocolo, sin cambios, da la misma respuesta cualitativa en el banco primario. Recortar
de dieciocho clases a ocho sube el F1 macro alineado de \num{0.5612} a \num{0.8052}, una mejora
aparente de \num{0.2440}, de la cual el control recupera \num{0.2155}, es decir el
\num{88.3}\,\%. En el punto de operación preregistrado, $K = 9$, el control recupera el
\num{87.1}\,\%, y el resultado se sostiene en los dos predictores que probamos.

Dos bancos, dos años, dos regiones, dos familias de modelo y un solo protocolo coinciden, por
tanto: manda el denominador. Donde difieren es en grado, y en la dirección que importa para quien
despliega: el banco de réplica, cuyo desbalance es un orden de magnitud peor, deja al mecanismo aún
menos que explicar.

\begin{figure}[tbp]
\centering
\includegraphics[width=\textwidth]{frontier-primary.pdf}
\caption{La misma frontera en el banco primario, para sus dos predictores más fuertes. El control
que no aplica mecanismo (cruces) ya recupera la mayor parte de la distancia recorrida por los
mecanismos, y la retirada por soporte (cuadrados) vuelve a quedar por debajo de los otros dos en
todo el recorrido.}
\label{fig:primario}
\end{figure}

\begin{figure}[tbp]
\centering
\includegraphics[width=\textwidth]{frontier-replication.pdf}
\caption{Frontera de calidad y cobertura en el banco de réplica, en los dos universos de clases
preregistrados. Las cuatro series son indistinguibles hasta $K = 6$ porque la cobertura es completa
y ningún mecanismo actúa; la línea punteada marca el criterio principal. La retirada por soporte
(cuadrados) es la única serie que baja conforme se retiran más clases.}
\label{fig:replica}
\end{figure}

\subsection{A igual cobertura, gana la abstención}
\label{sec:orden}

Una vez que la cobertura baja de uno y los mecanismos empiezan a diferir, se pueden ordenar. En el
criterio principal del banco de réplica, el rechazo por confianza alcanza \num{0.8303} frente a
\num{0.8055} del recorte de leyenda a la misma cobertura, una diferencia de \num{-0.0248} con
intervalo pareado al \num{95}\,\% de $(-0{,}0272,\, -0{,}0225)$ y $p$ corregido por Holm por debajo
de \num{0.0001}. Las dos regiones coinciden en signo y magnitud, con \num{-0.0248} y
\num{-0.0249}.

En el banco primario la diferencia tiene el mismo signo pero no se distingue de cero:
\num{-0.0288} con intervalo de $(-0{,}0414,\, +0{,}0080)$ y $p$ corregido por Holm de \num{0.336}.
El resumen honesto sobre los dos bancos es, por tanto, que no encontramos evidencia de que recortar
la leyenda supere a la abstención, y que en el banco mayor y más desbalanceado encontramos
evidencia significativa de lo contrario.

\begin{figure}[tbp]
\centering
\includegraphics[width=\textwidth]{cobertura.pdf}
\caption{Las mismas mediciones dibujadas contra la cobertura entregada en vez de contra el tamaño
de la leyenda, que es como las lee la literatura de clasificación selectiva. El control (cruces)
vive en la vertical de cobertura uno por construcción, porque nunca retiene nada; lo que la figura
enseña es cuántas de esas cruces quedan \emph{por encima} de las curvas de los mecanismos a
cobertura menor.}
\label{fig:cobertura}
\end{figure}

\subsection{La retirada por soporte es el criterio más débil, y no es monótona}
\label{sec:soporte}

Retirar clases por soporte bajo es el criterio que se reporta usar, y es el peor de los tres en los
dos bancos. En el banco de réplica alcanza \num{0.7129} en el criterio principal frente a
\num{0.8055} de la retirada por F1, una brecha de \num{0.0926} que se ensancha hasta \num{0.2085}
en $K = 3$. En el banco primario la brecha correspondiente está entre \num{0.12} y \num{0.13}.

Hace además algo que ninguno de los otros dos hace: \emph{empeora} conforme se retiran más clases,
de \num{0.7129} en $K = 5$ a \num{0.7041} en $K = 3$. Inspeccionar las leyendas lo explica. La
retirada por soporte quita la colza, que tiene \num{1718} parcelas y que el modelo separa bien,
mientras conserva las dos clases de pradera, que el modelo confunde entre sí. El criterio optimiza
el tamaño de muestra, que no es la misma cantidad que la calidad de lo que se entrega. La
figura~\ref{fig:leyendas} enseña la sustitución directamente.

\begin{figure}[tbp]
\centering
\includegraphics[width=\textwidth]{leyendas.pdf}
\caption{Qué clases conserva cada criterio conforme se recorta la leyenda, en el banco de réplica;
una celda llena significa que la clase se promete. La retirada por soporte (derecha) suelta la
colza en $K = 5$ mientras conserva las dos clases de pradera, que el modelo confunde entre sí; la
retirada por calidad (izquierda) hace lo contrario. Los dos criterios discrepan justamente en las
clases que deciden el resultado.}
\label{fig:leyendas}
\end{figure}

\subsection{Un meta-modelo entrenado retira clases de forma implícita}
\label{sec:arbitro}

El mecanismo que este artículo formaliza ya lo aplican, sin declararlo, los ensambles entrenados.
En el banco primario, un meta-modelo de apilamiento sobre tres miembros lleva una clase de
\num{355} parcelas a un F1 de exactamente \num{0.0000}, mientras que el mejor modelo individual
alcanza \num{0.7790} en esa misma clase. El meta-modelo ha retirado una clase del producto sin que
nadie lo decidiera, y sin que aparezca en ninguna tabla de resultados.

\begin{table}[tbp]
\caption{F1 por clase en el banco primario, libre de fuga, ordenado por soporte. El árbitro
entrenado es el único combinador que lleva una clase a cero, y lo hace con una de \num{355}
parcelas que todos los demás —incluida la media sin entrenar— manejan por encima de \num{0.77}.
La retirada es, por tanto, un comportamiento aprendido y no una propiedad de combinar.}
\label{tab:porclase}
\centering
\footnotesize
\setlength{\tabcolsep}{4.5pt}
\begin{tabular}{lrrrrr}
\toprule
Clase & soporte & mejor indiv. & media & voto pond. & árbitro \\
\midrule
Meadow & 6\,128 & 0,9099 & 0,8980 & 0,9046 & 0,9138 \\
Corn & 2\,196 & 0,9295 & 0,8683 & 0,9053 & 0,9274 \\
Grapevine & 2\,067 & 0,9066 & 0,9141 & 0,9154 & 0,9124 \\
Soft winter wheat & 1\,584 & 0,8922 & 0,8466 & 0,8712 & 0,8886 \\
Leguminous fodder & 660 & 0,5087 & 0,3498 & 0,4707 & 0,5243 \\
Winter rapeseed & 387 & 0,9358 & 0,9153 & 0,9275 & 0,9377 \\
\textbf{Winter durum wheat} & \textbf{355} & \textbf{0,7790} & \textbf{0,7752} & \textbf{0,7782} & \textbf{0,0000} \\
Spring barley & 198 & 0,6667 & 0,4691 & 0,5879 & 0,5228 \\
Mixed cereal & 193 & 0,4949 & 0,3525 & 0,4302 & 0,2709 \\
Potatoes & 103 & 0,5595 & 0,5143 & 0,5638 & 0,5823 \\
\bottomrule
\end{tabular}
\end{table}

\begin{table}[tbp]
\caption{El mismo meta-modelo de apilamiento bajo tres estimadores. No son tres opciones
plausibles: el primero está contaminado y el tercero lo castiga un artefacto de agregación.}
\label{tab:regimenes}
\centering
\small
\setlength{\tabcolsep}{5pt}
\begin{tabular}{lrl}
\toprule
Estimador & F1 macro & qué mide \\
\midrule
Reajuste sobre todas las parcelas & 0,7470 & dentro de muestra \\
Agrupado fuera de fold & 0,6789 & libre de fuga, puntuado una vez \\
Media de macros por bloque & 0,5360 & libre de fuga, sesgado al agregar \\
\bottomrule
\end{tabular}
\end{table}

El régimen en que se mide ese meta-modelo es decisivo. Puntuado dentro de muestra —reajustado sobre
todas las parcelas después de la validación cruzada, que es el comportamiento por omisión de muchas
implementaciones de apilamiento— reporta \num{0.7470}. Puntuado libre de fuga, agrupando sus
posteriores espaciales fuera de fold y puntuando una sola vez, reporta \num{0.6789}. Promediar los
macros por bloque en lugar de agrupar reporta \num{0.5360}. La fuga cuesta \num{0.068} y la forma
de agregar otros \num{0.143}; confundir las dos lleva a reportar \num{0.536} como la cifra honesta
cuando es \num{0.679}, y castiga al modelo por un artefacto de medición encima de su error real.
Frente al mejor modelo individual, el meta-modelo libre de fuga pierde por \num{0.0572} con un
intervalo de $(-0{,}0644,\, -0{,}0499)$ que excluye el cero.
